\documentclass[a4paper, 12pt]{article}

%%% SST LAB PROTOCOLL PREAMBLE
%%% 2019
%%%%%%%%%%%%%%%%%%%%%%%%%%%%%%%


%%% PACKAGES
%%%%%%%%%%%%%%%%%%%%%%%%%%%

\usepackage[ngerman]{babel}

\usepackage[utf8]{inputenc}
\usepackage{amsmath}
\usepackage{pgfplots}
\usepackage{tikz}
\usepackage[many]{tcolorbox}
\usepackage{graphicx}
\graphicspath{ {./graphics/} }
\usepackage{pdfpages}
\usepackage{dashrule}
\usepackage{float}
\usepackage{siunitx}
\usepackage{trfsigns}
\usepackage{booktabs}
\usepackage[european]{circuitikz}
\usepackage{tcolorbox}

%%% DOCUMENT GEOMETRY
%%%%%%%%%%%%%%%%%%%%%%%%%%%

\usepackage{geometry}
\geometry{
 a4paper,
 total={0.6180339887498948\paperwidth,0.6180339887498948\paperheight},
 top = 0.1458980337503154\paperheight,
 bottom = 0.1458980337503154\paperheight
 }
\setlength{\jot}{0.013155617496424828\paperheight}
\linespread{1.1458980337503154}

\setlength{\parskip}{0.013155617496424828\paperheight} % paragraph spacing


%%% COLORS
%%%%%%%%%%%%%%%%%%%%%%%%%%%

\definecolor{red1}{HTML}{f38181}
\definecolor{yellow1}{HTML}{fce38a}
\definecolor{green1}{HTML}{95e1d3}
\definecolor{blue1}{HTML}{66bfbf}
\definecolor{hsblue}{HTML}{00b1db}
\definecolor{hsgrey}{HTML}{afafaf}

%%% CONSTANTS
%%%%%%%%%%%%%%%%%%%%%%%%%%%
\newlength{\smallvert}
\setlength{\smallvert}{0.0131556\paperheight}


%%% COMMANDS
%%%%%%%%%%%%%%%%%%%%%%%%%%%

% differential d
\newcommand*\dif{\mathop{}\!\mathrm{d}}

% horizontal line
\newcommand{\holine}[1]{
  	\begin{center}
	  	\noindent{\color{hsgrey}\hdashrule[0ex]{#1}{1pt}{3mm}}\\%[0.0131556\paperheight]
  	\end{center}
}

% mini section
\newcommand{\minisec}[1]{ \noindent\underline{\textit {#1} } \\}

% quick function plot
\newcommand{\plotfun}[3]{
  \vspace{0.021286\paperheight}
  \begin{center}
    \begin{tikzpicture}
      \begin{axis}[
        axis x line=center,
        axis y line=center,
        ]
        \addplot[draw=red1][domain=#2:#3]{#1};
      \end{axis}
    \end{tikzpicture}
  \end{center}
}

% box for notes
\newcommand{\notebox}[1]{

\tcbset{colback=white,colframe=green1!100!black,title=Note!,width=0.618\paperwidth,arc=0pt}

 \begin{center}
  \begin{tcolorbox}[]
   #1 
  \end{tcolorbox}
 
 \end{center} 
 
}

% box for equation
\newcommand{\eqbox}[2]{
	
	\tcbset{colback=white,colframe=green1!100!black,title=,width=#2,arc=0pt}
	
	\begin{center}
		\begin{tcolorbox}[ams align*]
				#1
		\end{tcolorbox}
		
	\end{center} 
	
}
% END OF PREAMBLE

%%%%%%%%%%%%%%%%%%%%%%%%%%%%%%%%%%%%%

\begin{document}

%%%%%%%%%%%%%%%%%%%%%%%%%%%%%%%%%%%%%
  \includepdf{./titlepage/titlepage.pdf}
  \clearpage
  \setcounter{page}{1}
%%%%%%%%%%%%%%%%%%%%%%%%%%%%%%%%%%%%%

  \section{Register}
  Für die Zeitmessung wird die Capture Compare Unit 0 des Timer A mit dem zu
  messenden Eingangssignal an Pin \inlinecode{P2.2} verwendet. Zur Auswahl dieses Eingangs
  wird im \inlinecode{TACCTL0} \inlinecode{CCIS_1} gesetzt. Der Timer A wird
  außerdem in den continuous mode (\inlinecode{TACTL: MC_2}) gesetzt.

  Zum Zählen der Überläufe des Timers werden die Timerinterrupts aktiviert
  (\inlinecode{TAIE}). Zum Erfassen eines Capture-Events werden die Interrupts
  der CC-Einheit aktiviert (\inlinecode{TACCTL0: CCIE}).
  Um die Anzahl der Überläufe gering zu halten, wird der Taktvorteiler im
  \inlinecode{TACTL} Register auf 8 gesetzt (\inlinecode{TASSEL_1},
  \inlinecode{ID_3}). Jedoch muss beachtet werden, dass die Signalfrequenz
  nicht höher ist als die Zählfrequenz des Timers, da die Capturewerte sonst
  gleich sind (siehe Gleichung (1)).

  Zur Einstellung des Capture Modes muss das
  \inlinecode{CAP}-Flag im \inlinecode{TACCTL0} Register auf 1 und der Capture
  Mode \inlinecode{CM} des Captures je nach Messung gesetzt werden.
 
  \section{Ablauf der Messungen}

  \subsection{Periodenmessung}
  Bei der Periodenmessung wird der Capture Mode auf den Capture bei steigender
  Flanke gesetzt. Damit wird bei steigender Flanke des Eingangssignals die ISR
  des CCIFGs durchlaufen, in welcher dann beim ersten Durchlauf der erste
  Capture-Wert und beim zweiten der zweite gespeichert wird.

  \subsection{Impulsmessung}
  Bei der Impulsmessung wird der Capture zuerst bei steigender Flanke und dann
  bei fallender Flanke ausgelöst und die jeweiligen Capurewerte in der ISR aufgenommen.

  \subsection{Pausendauer}
  Die Pausendauermessung erfolgt analog der Impulsmessung nur dass zu Anfang der
  Capture Mode auf falling edge und in der ISR dann auf rising edge umgeschaltet wird.

  \section{Zeitberechnung}
  Die Gleichung für die Berechnung der Zeit ist für alle Messungen gleich.

  \begin{equation}
  T = \frac{\textrm{Capturewert2} - \textrm{Capturewert1}}{f_{CLK}}
\end{equation}

  Jedoch müssen die Timerüberläufe berücksichtigt werden, da die Capturewerte
  aus zwei verschiedenen Timerzählperioden stammen können und ohne
  Berücksichtigung der Überläufe nutzlos wären.

  \begin{equation}
    T = \frac{(N_{\textrm{Überläufe}} \cdot 2^{16}-1) + Z_2 - Z_1}{f_{CLK}}
  \end{equation}

  \section{Ausgabe einer Dezimalzahl}
  Ein 16 bit Integer kann mit der Laborfunktion \inlinecode{dis_zahl65535()} von
  Hr. Wenzel ausgegeben werden.

  Dabei bestimmt \inlinecode{csr} ob die Zahl an
  der momentanen LCD-Cursorposition (\inlinecode{csr=0}) oder an der durch die folgenden zwei
  Argumente angegebene Zeile und Spalte (\inlinecode{csr=1}) geschrieben werden soll. 
  \begin{lstlisting}
  int dis_zahl65535(unsigned char csr, 
                    unsigned char z, 
                    unsigned char s, 
                    unsigned int zahl); 
  \end{lstlisting}

\end{document}

% Local Variables:
% TeX-engine: luatex
% End:
